% ═══════════════════════════════════════════════════════════════
% LEHRERHINWEISE: Verbos en -ar (Regelmäßige -ar Verben)
% ═══════════════════════════════════════════════════════════════
% Sequenz 4: Mi colegio | Block c: Verbos en -ar
% ═══════════════════════════════════════════════════════════════

\documentclass[11pt, a4paper]{article}

% SW-MODUS TOGGLE
\newif\ifswmodus
\swmodusfalse

% PAKETE
\usepackage[utf8]{inputenc}
\usepackage[T1]{fontenc}
\usepackage[ngerman]{babel}
\usepackage{helvet}
\renewcommand{\familydefault}{\sfdefault}

\usepackage[margin=2cm]{geometry}
\usepackage{enumitem}
\usepackage{tabularx}
\usepackage{booktabs}
\usepackage{longtable}

\usepackage{xcolor}
\usepackage{amssymb}
\usepackage{tcolorbox}
\tcbuselibrary{skins}

\usepackage{fancyhdr}
\usepackage{lastpage}
\usepackage{needspace}

% FARBEN
\definecolor{spanisch-color}{HTML}{c41e3a}
\definecolor{grau-color}{HTML}{888888}
\definecolor{hellgrau-color}{HTML}{f5f5f5}
\definecolor{basis-color}{HTML}{2a7a4b}
\definecolor{standard-color}{HTML}{2c5aa0}
\definecolor{erweiterung-color}{HTML}{7b1fa2}
\definecolor{loesung-color}{HTML}{c62828}

\definecolor{sw-dark}{HTML}{000000}
\definecolor{sw-gray}{HTML}{666666}
\definecolor{sw-white}{HTML}{ffffff}

\ifswmodus
    \colorlet{spanisch}{sw-dark}
    \colorlet{grau}{sw-gray}
    \colorlet{hellgrau}{sw-white}
    \colorlet{basis}{sw-dark}
    \colorlet{standard}{sw-dark}
    \colorlet{erweiterung}{sw-dark}
    \colorlet{loesung}{sw-dark}
\else
    \colorlet{spanisch}{spanisch-color}
    \colorlet{grau}{grau-color}
    \colorlet{hellgrau}{hellgrau-color}
    \colorlet{basis}{basis-color}
    \colorlet{standard}{standard-color}
    \colorlet{erweiterung}{erweiterung-color}
    \colorlet{loesung}{loesung-color}
\fi

\newcommand{\fachfarbe}{spanisch}

% VARIABLEN
\newcommand{\thema}{Verbos en -ar}
\newcommand{\sequenz}{04}
\newcommand{\block}{c}
\newcommand{\fach}{Spanisch}
\newcommand{\klasse}{7}
\newcommand{\dauer}{90 Min. (Doppelstunde)}
\newcommand{\lerngruppengroesse}{24}
\newcommand{\datum}{\today}

% KOPF- UND FUSSZEILE
\pagestyle{fancy}
\fancyhf{}
\renewcommand{\headrulewidth}{0.4pt}
\renewcommand{\footrulewidth}{0.4pt}

\lhead{\fach{} -- Klasse \klasse}
\chead{\textbf{LH-\sequenz\block{} (Lehrerhinweise)}}
\rhead{\datum}

\lfoot{}
\cfoot{}
\rfoot{Seite \thepage\ von \pageref{LastPage}}

% BEFEHLE
\newcommand{\B}{\textcolor{basis}{\textbf{[B]}}}
\newcommand{\Std}{\textcolor{standard}{\textbf{[S]}}}
\newcommand{\E}{\textcolor{erweiterung}{\textbf{[E]}}}
\newcommand{\loesung}[1]{\textcolor{loesung}{\textbf{#1}}}
\newcommand{\es}[1]{\textcolor{\fachfarbe}{\textbf{#1}}}

\newtcolorbox{kurzplan}{
    colback=hellgrau,
    colframe=\fachfarbe,
    colbacktitle=white,
    coltitle=\fachfarbe,
    boxrule=1pt,
    arc=0pt,
    fonttitle=\bfseries,
    attach boxed title to top left={yshift=-2mm, xshift=5mm},
    boxed title style={boxrule=0pt, colback=white},
    title=Kurzplan
}

\newtcolorbox{differenzierung}{
    colback=white,
    colframe=grau,
    colbacktitle=white,
    coltitle=grau,
    boxrule=0.5pt,
    arc=0pt,
    fonttitle=\bfseries,
    attach boxed title to top left={yshift=-2mm, xshift=5mm},
    boxed title style={boxrule=0pt, colback=white},
    title=Differenzierung
}

\newtcolorbox{fehlerbox}{
    colback=white,
    colframe=loesung,
    colbacktitle=white,
    coltitle=loesung,
    boxrule=0.5pt,
    arc=0pt,
    fonttitle=\bfseries,
    attach boxed title to top left={yshift=-2mm, xshift=5mm},
    boxed title style={boxrule=0pt, colback=white},
    title=Häufige Fehler
}

\newtcolorbox{materialbox}{
    colback=white,
    colframe=\fachfarbe,
    colbacktitle=white,
    coltitle=\fachfarbe,
    boxrule=0.5pt,
    arc=0pt,
    fonttitle=\bfseries,
    attach boxed title to top left={yshift=-2mm, xshift=5mm},
    boxed title style={boxrule=0pt, colback=white},
    title=Material-Checkliste
}

% ═══════════════════════════════════════════════════════════════
% DOKUMENT
% ═══════════════════════════════════════════════════════════════

\begin{document}

% KOPFBEREICH
\begin{center}
    {\Large\bfseries\textcolor{\fachfarbe}{Lehrerhinweise: \thema}}\\[0.3em]
    {\normalsize LH-\sequenz\block{} | \dauer}
\end{center}

\vspace{10pt}

% ═══════════════════════════════════════════════════════════════
% 1. KURZPLAN
% ═══════════════════════════════════════════════════════════════

\section*{1. Stundenverlauf (Kurzplan)}

\textbf{1. Stunde (45 Min.)}

\begin{tabularx}{\textwidth}{|c|l|X|l|}
\hline
\textbf{Zeit} & \textbf{Phase} & \textbf{Inhalt} & \textbf{Material} \\
\hline
0--5 & Einstieg & Bilder: Aktivitäten (sprechen, singen, tanzen) & Bilder \\
\hline
5--15 & Input & \es{hablar} als Modellverb, Konjugation an Tafel & Tafel \\
\hline
15--25 & Übung 1 & Chorisches Sprechen: ,,hablo, hablas, habla, hablamos'' & -- \\
\hline
25--35 & Übung 2 & AB Aufgabe 1: Endungen einsetzen & AB \\
\hline
35--45 & Sicherung & Ergebnisse vergleichen, Fehler korrigieren & Tafel \\
\hline
\end{tabularx}

\vspace{1em}

\textbf{2. Stunde (45 Min.)}

\begin{tabularx}{\textwidth}{|c|l|X|l|}
\hline
\textbf{Zeit} & \textbf{Phase} & \textbf{Inhalt} & \textbf{Material} \\
\hline
0--5 & Wiederholung & Schnelles Endungen-Quiz mündlich & -- \\
\hline
5--15 & Input & Weitere Verben: estudiar, escuchar, trabajar, cantar, bailar & Tafel \\
\hline
15--30 & Anwendung & AB Aufgaben 2-3: Konjugieren, Sätze bilden & AB \\
\hline
30--40 & Festigung & Partner-Interview: ,,¿Qué estudias?'' & -- \\
\hline
40--45 & Abschluss & Zusammenfassung, HA: AB Aufgabe 4 & Tafel \\
\hline
\end{tabularx}

% ═══════════════════════════════════════════════════════════════
% 2. DIFFERENZIERUNG
% ═══════════════════════════════════════════════════════════════

\section*{2. Differenzierung}

\begin{differenzierung}
\textbf{Legende:} \B{} Basis (BR) | \Std{} Standard (MR) | \E{} Erweiterung (GY)

\vspace{0.5em}

\textbf{WICHTIG:} Diese Marker sind NUR für die Lehrkraft. Auf Schülermaterial erscheinen KEINE Niveaumarkierungen!
\end{differenzierung}

\vspace{1em}

\begin{tabularx}{\textwidth}{|l|X|X|X|}
\hline
& \B{} \textbf{Basis} & \Std{} \textbf{Standard} & \E{} \textbf{Erweiterung} \\
\hline
\textbf{Personen} & yo, tú, él/ella & + nosotros & + vosotros, ellos \\
\hline
\textbf{Aufg. 1} & Mit Konjugationstabelle & Selbstständig & Ohne Hilfsmittel \\
\hline
\textbf{Aufg. 2} & Mit Beispielen & Selbstständig & + Erklärung \\
\hline
\textbf{Aufg. 3} & 4 Sätze & 6 Sätze & 6 Sätze + Variation \\
\hline
\textbf{Aufg. 4} & 3 einfache Sätze & 4 Sätze mit Vielfalt & 5+ Sätze, komplex \\
\hline
\end{tabularx}

\vspace{1em}

\textbf{Konkrete Hinweise:}
\begin{itemize}
    \item \B{} SuS mit LRS: Konjugationstabelle als Hilfskarte, Endungen farbig markiert
    \item \B{} DaZ-SuS: Parallele deutsche Konjugation zeigen (ich spreche, du sprichst...)
    \item \E{} Leistungsstarke: Alle 6 Personen lernen, zusätzliche Verben (preguntar, contestar)
\end{itemize}

% ═══════════════════════════════════════════════════════════════
% 3. HÄUFIGE FEHLER
% ═══════════════════════════════════════════════════════════════

\section*{3. Häufige Fehler}

\begin{fehlerbox}
\begin{tabularx}{\textwidth}{|l|X|X|}
\hline
\textbf{Fehler} & \textbf{Beispiel} & \textbf{Reaktion} \\
\hline
-as / -a verwechselt & *,,Tú habla'' statt \es{Tú hablas} & Personalpronomen betonen, farbliche Zuordnung \\
\hline
Stamm falsch & *,,Yo canta'' statt \es{Yo canto} & Regel: Stamm bleibt gleich, nur Endung ändert sich \\
\hline
Akzent vergessen & *,,hablamos'' mit falschem Akzent & Betonung üben: habl-A-mos \\
\hline
-amos vs. -emos & *,,hablemos'' statt \es{hablamos} & -amos nur bei -ar Verben! \\
\hline
Subjekt vergessen & Wenn ohne Kontext & Vorerst mit Pronomen üben \\
\hline
\end{tabularx}
\end{fehlerbox}

% ═══════════════════════════════════════════════════════════════
% 4. MATERIAL-CHECKLISTE
% ═══════════════════════════════════════════════════════════════

\section*{4. Material-Checkliste}

\begin{materialbox}
\begin{itemize}[label=$\square$]
    \item AB-\sequenz\block.pdf (\lerngruppengroesse{} Kopien)
    \item ML-\sequenz\block.pdf (1$\times$ für Lehrkraft)
    \item Lehrwerk: Qué pasa 1, S. 28-30
    \item Tafelmarker / Kreide (verschiedene Farben für Endungen)
    \item Optional: LP-\sequenz\block.html (Tablets/PCs)
    \item Optional: Bildkarten mit Aktivitäten (sprechen, singen, tanzen, etc.)
    \item Optional: Konjugationshilfe-Karten für \B{}-SuS
\end{itemize}
\end{materialbox}

% ═══════════════════════════════════════════════════════════════
% 5. LÖSUNGEN
% ═══════════════════════════════════════════════════════════════

\section*{5. Lösungen AB-\sequenz\block}

\textbf{Merkkasten: Endungen und Verben}

\begin{tabular}{ll}
Endungen: & -o, -as, -a, -amos \\[0.3em]
hablar: & hablo, hablas, habla, hablamos \\[0.3em]
\end{tabular}

\begin{tabular}{ll@{\hspace{2cm}}ll}
hablar & = sprechen & trabajar & = arbeiten \\
estudiar & = lernen, studieren & cantar & = singen \\
escuchar & = hören, zuhören & bailar & = tanzen \\
\end{tabular}

\vspace{1em}

\textbf{Aufgabe 1: Endungen einsetzen} (4 Punkte)

\begin{tabular}{ll}
a) & habl\loesung{o} \\
b) & escuch\loesung{as} \\
c) & estudi\loesung{amos} \\
d) & trabaj\loesung{a} \\
\end{tabular}

\vspace{1em}

\textbf{Aufgabe 2: Verben konjugieren} (4 Punkte)

\begin{tabular}{ll}
a) & \loesung{canto} \\
b) & \loesung{bailas} \\
c) & \loesung{estudia} \\
d) & \loesung{trabajamos} \\
\end{tabular}

\vspace{1em}

\textbf{Aufgabe 3: Sätze bilden} (6 Punkte)

\begin{tabular}{ll}
a) & \loesung{Yo hablo español.} \\
b) & \loesung{Tú escuchas al profesor.} \\
c) & \loesung{María estudia historia.} \\
d) & \loesung{Nosotros cantamos una canción.} \\
e) & \loesung{Pablo trabaja en clase.} \\
f) & \loesung{Tú bailas muy bien.} \\
\end{tabular}

\vspace{1em}

\textbf{Aufgabe 4: Schulalltag beschreiben} (6 Punkte)

Bewertungskriterien:
\begin{itemize}
    \item 4 verschiedene -ar Verben verwendet: 2P
    \item Korrekte Konjugation: 2P
    \item Sinnvoller Inhalt (Schulbezug): 1P
    \item Sprachliche Korrektheit: 1P
\end{itemize}

\textit{Beispiellösungen:}
\begin{itemize}
    \item \loesung{Yo estudio español en el colegio.}
    \item \loesung{En clase escucho al profesor.}
    \item \loesung{Nosotros trabajamos mucho en matemáticas.}
    \item \loesung{En música canto con mis amigos.}
\end{itemize}

% ═══════════════════════════════════════════════════════════════
% 6. HAUSAUFGABE
% ═══════════════════════════════════════════════════════════════

\section*{6. Hausaufgabe}

\begin{itemize}
    \item AB-04c Aufgabe 4 fertigstellen (falls nicht in der Stunde geschafft)
    \item Vokabeln lernen: hablar, estudiar, escuchar, trabajar, cantar, bailar
    \item Konjugationsendungen auswendig: -o, -as, -a, -amos (-áis, -an für \E{})
    \item Optional: Lernpfad LP-\sequenz\block.html bearbeiten
\end{itemize}

% ═══════════════════════════════════════════════════════════════
% 7. REFLEXIONSNOTIZEN
% ═══════════════════════════════════════════════════════════════

\section*{7. Reflexionsnotizen (nach Durchführung)}

\begin{tabularx}{\textwidth}{|l|X|}
\hline
\textbf{Was lief gut?} & \\[2em]
\hline
\textbf{Was lief nicht?} & \\[2em]
\hline
\textbf{Zeitmanagement} & \\[2em]
\hline
\textbf{Für nächstes Mal} & \\[2em]
\hline
\end{tabularx}

\end{document}
